\documentclass[11pt,a4paper]{article}
\usepackage[utf8]{inputenc}
\usepackage[T1]{fontenc}
\usepackage{amsmath,amssymb,amsthm}
\usepackage{graphicx}
\usepackage{hyperref}
\usepackage{booktabs}
\usepackage{algorithm}
\usepackage{algpseudocode}
\usepackage{geometry}
\usepackage{natbib}
\usepackage{authblk}
\geometry{margin=1in}

\title{\textbf{Thermodynamic Epistemic Routing: \\
Zero-Parameter Detection of Factual vs. Speculative Generation in Large Language Models}}

\author{Scott Seto}
\affil{Independent Researcher \\ \texttt{https://github.com/ameritusweb/thermodynamic-epistemic-routing}}

\date{\today}

\begin{document}

\maketitle

\begin{abstract}
Large language models (LLMs) exhibit remarkable generative capabilities but suffer from a fundamental reliability issue: they cannot reliably distinguish between factual recall and speculative extrapolation. We present a novel approach that induces models to develop thermodynamically distinct computational pathways for these two epistemic modes. Unlike prior work that attempts to classify activation patterns spatially, we discover that the \emph{dynamics} of layer-to-layer hidden state changes provide a more fundamental signal. Through adversarial co-training with LoRA adapters and a predictor network, we demonstrate that delta magnitude patterns---the L2 norms of consecutive hidden state differences---naturally separate factual and speculative content. This thermodynamic signal enables \textbf{98.39\% classification accuracy on Qwen2.5-1.5B using only a scalar threshold with zero additional parameters} (gap = 13.62 units), compared to 91.15\% for the best spatial MLP probe. \textbf{Direct replication on Qwen2.5-7B (4.67$\times$ parameters, identical hyperparameters) yields 98.13\% accuracy with gap = 13.88}---confirming the method transfers across model scale without modification. Our key finding is that spatial activation patterns read a downstream \emph{echo} of the true signal, while thermodynamic features access the primary output of epistemic specialization. The approach preserves generation quality (eval loss improved from 1.914 to 1.835 at 1.5B; 2.323 to 1.076 at 7B) while creating interpretable epistemic structure, offering a path toward mechanistically grounded uncertainty quantification in production LLM systems.

\textbf{Keywords:} Large Language Models, Epistemic Uncertainty, Mechanistic Interpretability, Adversarial Training, LoRA Fine-tuning, Hallucination Detection
\end{abstract}

\section{Introduction}

\subsection{The Epistemic Honesty Problem}

Large language models have achieved remarkable performance on diverse tasks \citep{brown2020language, touvron2023llama, jiang2023mistral}, yet they systematically fail to distinguish between factual recall and speculative extrapolation. A model answering ``What is the capital of France?'' (factual) exhibits no detectable internal difference from one answering ``What might future cities on Mars look like?'' (speculative). Both are generated with equal confidence, making post-hoc hallucination detection challenging.

Current mitigation strategies fall into several categories: retrieval-augmented generation (RAG) \citep{lewis2020retrieval}, post-hoc fact-checking against knowledge bases, constitutional AI training to encourage refusals \citep{bai2022constitutional}, and uncertainty quantification via sampling \citep{kuhn2023semantic}. While valuable, these approaches share key limitations: (1) they treat hallucination as an output-level problem rather than an internal representational issue, (2) they often reduce model capability through conservative training, (3) they lack mechanistic interpretability, and (4) they provide no pre-generation detection mechanism.

\subsection{Our Approach: Thermodynamic Routing}

We reframe the problem: rather than asking ``how do we stop the model from guessing?'', we ask ``how do we make the model's epistemic state legible in its internal dynamics?'' Our key insight is that different epistemic processes should leave distinct \emph{thermodynamic signatures} in the forward pass---measurable properties of how hidden states evolve layer-to-layer, independent of what specific content they represent.

We demonstrate this through adversarial co-training: a generator (base LLM with LoRA adapters) and a predictor (lightweight MLP) evolve together, with the generator receiving gradient pressure to develop thermodynamically distinct pathways for factual versus speculative content. The training signal leverages the \emph{context window boundary as a ground truth oracle}: questions answerable from context are factual, those requiring extrapolation are speculative.

\subsection{Key Contributions}

\begin{enumerate}
\item \textbf{Discovery of thermodynamic routing signal}: We identify delta magnitudes $T(x) = \text{mean}(\|h_{n+1} - h_n\|)$ as a superior epistemic signal compared to spatial activation patterns, achieving 98.39\% classification accuracy with zero learned parameters.

\item \textbf{Explanation of spatial probe ceiling}: We demonstrate that spatial MLP probes (91.15\% accuracy) read a downstream \emph{echo} of the thermodynamic signal, not its source, explaining the 7pp performance gap.

\item \textbf{Evidence that specialization improves generation}: Contrary to typical capability-safety tradeoffs, thermodynamic specialization \emph{improved} generation quality (eval loss 1.914 $\rightarrow$ 1.835), suggesting uniform computation was a suboptimal local equilibrium.

\item \textbf{Scalable training methodology}: The approach requires no architectural changes, uses automated labeling (no human annotation), and produces interpretable epistemic structure suitable for production deployment.
\end{enumerate}

\section{Related Work}

\subsection{Uncertainty Quantification in Neural Networks}

Prior work on uncertainty in neural networks includes Bayesian approaches \citep{blundell2015weight, gal2016dropout}, ensemble methods \citep{lakshminarayanan2017simple}, and conformal prediction \citep{angelopoulos2021gentle}. For LLMs specifically, recent work has explored semantic entropy \citep{kuhn2023semantic}, self-consistency \citep{wang2023selfconsistency}, and verbalized confidence \citep{tian2023finetuning}. Our approach differs by operating on \emph{internal dynamics} rather than output distributions, creating structural changes during training rather than post-hoc estimation, and requiring only a single forward pass.

\subsection{Mechanistic Interpretability}

Recent interpretability research includes probing classifiers \citep{belinkov2021probing}, activation patching \citep{meng2022locating}, and sparse autoencoders \citep{cunningham2023sparse}. Our work contributes by using interpretability tools (probes) as part of the \emph{training objective} itself, creating interpretable structure through gradient pressure rather than discovering it post-hoc. Critically, we identify that spatial probes can converge to spurious correlates, and thermodynamic features provide more fundamental access to epistemic state.

\subsection{Geometric and Trajectory-Based Interpretability}

Recent work has begun treating the transformer residual stream as a geometric trajectory through high-dimensional space rather than a sequence of independent layer outputs. \citet{manson2025curved} introduce \emph{Curved Inference}, which measures two trajectory properties under a pullback metric $G = U^\top U$ derived from the unembedding matrix $U$: \emph{salience} $S(t) = \|x_{t+1} - x_t\|_G$ (layer-to-layer movement weighted by token probability sensitivity) and \emph{curvature} $\kappa_i$ (sharpness of directional change). Applied to Gemma3-1b and LLaMA3.2-3b across semantically shifted prompts, the study finds that concern-shifted inputs reliably alter trajectory geometry, and identifies a strong anticorrelation between salience and curvature ($r = -0.89$ in LLaMA)---interpreted as a \emph{representational budget}: the model either moves far or reorients sharply, not both simultaneously.

Our work converges independently on the same core intuition: that layer-to-layer hidden state movement, rather than activation geometry at any single layer, carries the most fundamental epistemic signal. The principal methodological differences are threefold. First, Manson's salience is computed under the semantically grounded pullback metric while our $T(x) = \text{mean}(\|h_{n+1} - h_n\|)$ uses the Euclidean norm---geometrically unweighted but computationally free. Second, Curved Inference is a passive, inference-only analysis; our approach \emph{actively reshapes} dynamics through LoRA training, forcing the gap open rather than measuring a pre-existing one. Third, validation scale differs by three orders of magnitude: 20 hand-curated prompts versus 20,000 behaviorally labeled examples with ground-truth epistemic annotations. The two papers are thus complementary: geometric theory first in Manson, behavioral validation at scale here.

\subsection{Adversarial and Co-evolutionary Training}

GANs \citep{goodfellow2014generative} pioneered adversarial training for generative models. More recently, debate-based approaches have been proposed for AI alignment \citep{irving2018ai}. Our architecture shares the two-model competitive structure but differs critically: the generator is not trying to fool the predictor but to honestly route content, the predictor has perfect ground truth during training, and the goal is internal structural change rather than output quality.

\subsection{Parameter-Efficient Fine-Tuning}

LoRA \citep{hu2021lora} and related methods \citep{lialin2023scaling} enable efficient adaptation of large models. We leverage LoRA not for task-specific adaptation but to reshape internal dynamics while preserving base model capabilities.

\section{Methodology}

\subsection{Problem Formulation}

Let $G$ be a transformer-based language model with hidden dimension $d$ and $L$ layers. For an input sequence $x$, let $h_\ell \in \mathbb{R}^d$ denote the hidden state at layer $\ell$ and position $t$ (last valid token). We define:

\begin{equation}
\text{Epistemic state}(x) = \begin{cases}
\text{factual} & \text{if } \text{answer}(x) \in \text{context}(x) \\
\text{speculative} & \text{otherwise}
\end{cases}
\end{equation}

The goal is to train $G$ such that its internal representations $\{h_\ell\}_{\ell=1}^L$ exhibit detectable, consistent differences based on epistemic state, enabling pre-generation classification.

\subsection{Training Signal: Context Window as Oracle}

We construct a dataset $\mathcal{D} = \{(c_i, q_i^f, a_i^f, q_i^s, a_i^s)\}$ where:
\begin{itemize}
\item $c_i$ is a context passage (e.g., Wikipedia paragraph)
\item $q_i^f$ is a factual question whose answer $a_i^f$ appears verbatim in $c_i$
\item $q_i^s$ is a speculative question whose answer $a_i^s$ requires extrapolation beyond $c_i$
\end{itemize}

Ground truth labels are \emph{automated}: factual if answerable from context, speculative otherwise. To prevent phrasing artifacts, question pairs are constrained to share identical grammatical structure (same question word, sentence pattern).

\subsection{Architecture}

Our system comprises two components:

\textbf{Generator $G$:} Base LLM (Qwen2.5-1.5B) augmented with LoRA adapters \citep{hu2021lora}. LoRA parameters $\theta_\text{LoRA}$ modify projection matrices in all transformer layers with rank $r=128$, alpha $\alpha=256$.

\textbf{Predictor $P$:} Lightweight MLP classifier observing activation patterns from $G$. Architecture: $[\text{input\_dim} \rightarrow 256 \rightarrow 128 \rightarrow 1]$ with LayerNorm, ReLU, dropout $[0.5, 0.2]$, trained with weighted binary cross-entropy.

\subsection{The Search for a Programmable Signal}

We explored multiple candidate signals before identifying the thermodynamic approach:

\textbf{Attempt 1 --- Spatial Separation:} Centroid L2 distance $\|c_\text{fact} - c_\text{spec}\|_2$ at target layer. \emph{Result:} Programmable, achieves 91\% but fails to scale (single-endpoint constraint).

\textbf{Attempt 2 --- Raw Turbulence:} Cosine distance $1 - \cos(h_n, h_{n+1})$ between consecutive states. \emph{Result:} Consistent diagnostic signal but gradient vanishes due to near-parallelism of residual stream vectors ($h_n$ and $h_{n+1}$ differ by small perturbations).

\textbf{Attempt 3 --- Delta Curvature:} $\cos(\delta_n, \delta_{n+1})$ where $\delta_n = h_{n+1} - h_n$. \emph{Result:} Fixes gradient geometry but no consistent directional signal (deltas are nearly orthogonal).

\textbf{Attempt 4 --- Delta Magnitude (Success):} $T(x) = \text{mean}(\|\delta_n\|_2)$ across target layers. \emph{Result:} Gradient never vanishes ($\nabla_h \|\delta\| = \delta / \|\delta\|$), empirically separates classes, sign always positive.

\subsection{ThermoSpatial Loss Function}

The final training objective combines thermodynamic and spatial constraints:

\begin{align}
\mathcal{L}_\text{thermo} &= \text{mean}\left(\text{ReLU}\left(m_\text{thermo} - (T_\text{fact}^{(i)} - T_\text{spec}^{(i)})\right)\right) \label{eq:thermo} \\
\mathcal{L}_\text{spatial} &= \text{ReLU}\left(m_\text{spatial} - \|c_\text{fact} - c_\text{spec}\|_2\right) \label{eq:spatial} \\
\mathcal{L}_\text{routing} &= \mathcal{L}_\text{BCE} + \lambda_\text{thermo} \mathcal{L}_\text{thermo} + \lambda_\text{contrastive} \mathcal{L}_\text{spatial} \label{eq:routing} \\
\mathcal{L}_\text{total} &= \mathcal{L}_\text{generation} + \lambda_\text{routing} \mathcal{L}_\text{routing} \label{eq:total}
\end{align}

Key design decisions:
\begin{itemize}
\item \textbf{Pairwise, context-normalized}: $T_\text{fact}^{(i)} - T_\text{spec}^{(i)}$ computed per same-context pair controls for context complexity.
\item \textbf{Scale calibration}: Delta magnitude losses are $O(\|h\|) \approx 25$, not $O(1)$. We use $\lambda_\text{thermo} = 0.1$ (not 5.0) to keep thermo contribution $\sim 0.2$, balanced against generation loss $\sim 1.5$.
\item \textbf{Margin at natural gap}: $m_\text{thermo} = 4.0$ set at baseline gap, making loss a floor constraint (fires only when separation degrades), not a ceiling.
\end{itemize}

\subsection{Adversarial Co-Evolution Loop}

Training alternates between two update types:

\textbf{Generator update (every step):}
\begin{itemize}
\item Predictor in eval mode, optimizer disabled
\item Routing loss gradient flows: $(\mathcal{L}_\text{BCE} + \mathcal{L}_\text{thermo} + \mathcal{L}_\text{spatial}) \rightarrow \{h_\ell\} \rightarrow \theta_\text{LoRA}$
\item LoRA learns to produce activations satisfying thermodynamic constraint
\end{itemize}

\textbf{Predictor update (every $N$ steps after warmup):}
\begin{itemize}
\item Activations detached from computation graph
\item Only predictor parameters receive gradient
\item Predictor adapts to shifting activation space as LoRA evolves
\end{itemize}

The co-evolution creates mutual pressure: as generator produces better-separated activations, predictor must work harder; as predictor improves, generator must widen gap further.

\section{Experimental Setup}

\subsection{Dataset}

\textbf{Source:} SQuAD v2 \citep{rajpurkar2018know} Wikipedia contexts ($\sim$10,000 passages).

\textbf{Generation:} Claude Haiku (claude-haiku-4-5-20251001) via Anthropic Batch API.

\textbf{Format:} Each context yields one factual and one speculative question-answer pair, constrained to identical grammatical structure.

\textbf{Cleaning:} Removed 1,386 refusal-leakage examples (``The context does not state...'').

\textbf{Augmentation:} Added 1,991 targeted examples (991 hard-negatives: confident but ungrounded; 1,000 complex-factual: multi-hop reasoning).

\textbf{Final split:} 16,889 train (8,891 factual / 7,998 speculative), 1,845 validation, 1,869 test.

\subsection{Model Configuration}

\textbf{Base model:} Qwen2.5-1.5B-Instruct \citep{qwen2024}, $d=1536$, $L=28$ layers.

\textbf{LoRA:} $r=128$, $\alpha=256$, dropout $0.1$, all projection matrices (q, k, v, o, gate, up, down).

\textbf{Thermodynamic layers:} $\{16, 17, 18, 19\}$ (within the ``blushing zone'' identified by layer sweep).

\textbf{Training:} 6 epochs, batch size 4 $\times$ 32 gradient accumulation (effective 128), AdamW, gradient checkpointing enabled.

\textbf{Hyperparameters:} $\lambda_\text{routing} = 1.0$, $\lambda_\text{thermo} = 0.1$, $\lambda_\text{contrastive} = 1.0$, $m_\text{thermo} = 4.0$, $m_\text{spatial} = 4.0$.

\subsection{Evaluation Metrics}

\textbf{Primary metric:} Classification accuracy on held-out test set (1,869 examples).

\textbf{Thermodynamic signal:} Mean factual $T(x)$, mean speculative $T(x)$, gap, standard deviations.

\textbf{Error modes:} FM1 (factual $\rightarrow$ speculative), FM2 (speculative $\rightarrow$ factual).

\textbf{Generation quality:} Eval language modeling loss on standard prompts.

\textbf{Topology:} PCA silhouette score, Davies-Bouldin index, separation ratio.

\section{Results}

\subsection{Layer Sweep: The ``Blushing Zone''}

Before training, we probed all 28 layers with linear classifiers (logistic regression, 2,000 examples). Epistemic signal peaks at layers 17--20 with 83--84\% accuracy, significantly exceeding the penultimate layer 26 (74\%) previously assumed optimal.

\textbf{Post-LoRA:} Blushing zone layers improve to 88--89\%, confirming routing gradient propagates through the full network.

\textbf{Implication:} Mid-network representations carry stronger epistemic signal than near-output layers where signal dilutes before vocabulary projection.

\subsection{Spatial Probe Performance Progression}

Table~\ref{tab:spatial} shows spatial MLP probe accuracy across experimental phases.

\begin{table}[h]
\centering
\caption{Spatial MLP predictor accuracy progression (test set).}
\label{tab:spatial}
\begin{tabular}{lccc}
\toprule
Method & Overall & Factual & Speculative \\
\midrule
Phase 1: Frozen baseline, layer 26 & 85.9\% & $\sim$85\% & $\sim$87\% \\
Phase 3 Run 2: LoRA (r=128, 6ep), layer 26 & 90.45\% & 93.8\% & 87.0\% \\
Multi-feature: Layers 17--19, 5-token pool & \textbf{91.15\%} & 93.7\% & 88.6\% \\
Phase 6 MLP: Thermo LoRA, layer 26 & 89.51\% & 91.51\% & 87.15\% \\
\bottomrule
\end{tabular}
\end{table}

\textbf{Finding:} Spatial probes plateau at $\sim$91\% despite improved LoRA training, larger feature dimensions (4,608-dim from 3 layers), and augmented data. This ceiling persists across configurations, suggesting a fundamental limitation.

\subsection{Thermodynamic Signal Emergence}

Table~\ref{tab:thermo_progression} shows thermodynamic signal evolution during Phase 6 training.

\begin{table}[h]
\centering
\caption{Thermodynamic signal progression across epochs (delta magnitudes). Bottom rows show 7B epoch 6 final for direct comparison.}
\label{tab:thermo_progression}
\begin{tabular}{lcccccc}
\toprule
Metric & Baseline & Epoch 1 & Epoch 2 & Epoch 3 & Epoch 6 \\
\midrule
\multicolumn{6}{l}{\textit{Qwen2.5-1.5B (Phase 6)}} \\
Factual $T(x)$ & $\sim$26 & $\sim$27 & $\sim$30 & $\sim$32 & \textbf{36.22} \\
Speculative $T(x)$ & $\sim$22.5 & $\sim$20.5 & $\sim$21 & $\sim$22 & \textbf{22.59} \\
Gap & $\sim$3.5 & $\sim$5.5 & $\sim$8.5 & $\sim$10 & \textbf{13.62} \\
$\mathcal{L}_\text{thermo}$ & --- & 1.0--2.5 & 0--0.5 & $\sim$0 & $\sim$0 \\
Eval LM loss & --- & 1.914 & --- & 1.842 & \textbf{1.835} \\
\midrule
\multicolumn{6}{l}{\textit{Qwen2.5-7B (Phase 8, identical hyperparameters, epoch 6 final)}} \\
Factual $T(x)$ & 4.695 (natural) & --- & --- & --- & \textbf{43.96} \\
Speculative $T(x)$ & --- & --- & --- & --- & \textbf{30.08} \\
Gap & $\sim$4.7 & --- & --- & --- & \textbf{13.88} \\
Eval LM loss & --- & 2.323 & --- & --- & \textbf{1.076} \\
\bottomrule
\end{tabular}
\end{table}

\textbf{Key observations:}

\begin{enumerate}
\item \textbf{Bidirectional specialization} (Epoch 2): Model spontaneously pushes factual up and speculative down simultaneously, discovering that bidirectional movement doubles gradient signal.

\item \textbf{Constraint goes slack} (Epoch 3): $\mathcal{L}_\text{thermo} \approx 0$ once gap exceeds $m_\text{thermo} = 4.0$. Gap continues widening in subsequent epochs \emph{without loss pressure}, indicating internalized specialization.

\item \textbf{Generation improves with specialization}: Eval loss decreases monotonically from 1.914 to 1.835 as gap widens. Thermodynamic specialization is not a capability-safety tradeoff but an \emph{optimal generation strategy} the model discovered.
\end{enumerate}

\subsection{Zero-Parameter Threshold Evaluation}

With gap = 13.62 (3.4$\times$ the required margin), we evaluate a scalar threshold $\tau$ on $T(x)$:

\begin{equation}
\text{classify}(x) = \begin{cases}
\text{factual} & \text{if } T(x) > \tau \\
\text{speculative} & \text{otherwise}
\end{cases}
\end{equation}

Table~\ref{tab:threshold} shows performance across threshold values.

\begin{table}[h]
\centering
\caption{Threshold sweep results (test set, 1,869 examples). Selected thresholds shown; best highlighted per model.}
\label{tab:threshold}
\begin{tabular}{lrcccccc}
\toprule
Model & $\tau$ & Overall & Factual & Speculative & FM1 & FM2 \\
\midrule
\multicolumn{7}{l}{\textit{Qwen2.5-1.5B} \ (factual $\mu$=36.22, spec $\mu$=22.59, gap=13.62)} \\
& 25.0 & 96.04\% & 99.90\% & 91.47\% & 1 & 73 \\
& 26.0 & 97.70\% & 99.70\% & 95.33\% & 3 & 40 \\
& \textbf{27.2} & \textbf{98.39\%} & \textbf{99.31\%} & \textbf{97.31\%} & \textbf{7} & \textbf{23} \\
& 28.0 & 98.23\% & 98.72\% & 97.66\% & 13 & 20 \\
& 29.0 & 97.91\% & 97.93\% & 97.90\% & 21 & 18 \\
\midrule
\multicolumn{7}{l}{\textit{Qwen2.5-7B} \ (factual $\mu$=43.96, spec $\mu$=30.08, gap=13.88)} \\
& 33.8 & 97.38\% & 99.41\% & 94.98\% & 6 & 43 \\
& 34.8 & 97.81\% & 98.91\% & 96.50\% & 11 & 30 \\
& \textbf{35.4} & \textbf{98.13\%} & \textbf{98.72\%} & \textbf{97.43\%} & \textbf{13} & \textbf{22} \\
& 35.9 & 97.91\% & 98.22\% & 97.55\% & 18 & 21 \\
& 36.7 & 97.43\% & 97.24\% & 97.66\% & 28 & 20 \\
\bottomrule
\end{tabular}
\end{table}

\textbf{Optimal thresholds:}
\begin{itemize}
\item \textbf{1.5B} ($\tau = 27.2$): Overall \textbf{98.39\%}, factual 99.31\%, speculative 97.31\% (FM1 0.7\%, FM2 2.7\%)
\item \textbf{7B} ($\tau = 35.4$): Overall \textbf{98.13\%}, factual 98.72\%, speculative 97.43\% (FM1 1.3\%, FM2 2.6\%)
\item Additional parameters: \textbf{0} (both models)
\end{itemize}

\textbf{Distribution statistics:}
\begin{itemize}
\item 1.5B --- Factual: mean 36.22, std 3.42 \ | \ Speculative: mean 22.59, std 2.24 \ | \ Gap: 13.62 (6.1 pooled std)
\item 7B \ \ \ --- Factual: mean 43.96, std 3.65 \ | \ Speculative: mean 30.08, std 2.27 \ | \ Gap: 13.88 (5.9 pooled std)
\end{itemize}

\subsection{Comparison: Spatial MLP vs. Zero-Parameter Threshold}

Table~\ref{tab:comparison} compares the best spatial probe with the thermodynamic threshold on the \emph{same Phase 6 model}.

\begin{table}[h]
\centering
\caption{Spatial probe vs. thermodynamic threshold. 7B replication uses identical hyperparameters; threshold recalibrated on held-out set.}
\label{tab:comparison}
\begin{tabular}{llccc}
\toprule
Model & Method & Signal & Parameters & Accuracy \\
\midrule
1.5B & Spatial MLP (layer 26) & $h_{26}$ & $\sim$1.2M & 89.51\% \\
1.5B & \textbf{T(x) threshold} & $\|\delta_n\|$ & \textbf{0} & \textbf{98.39\%} \\
\midrule
7B & \textbf{T(x) threshold} & $\|\delta_n\|$ & \textbf{0} & \textbf{98.13\%} \\
\bottomrule
\end{tabular}
\end{table}

\textbf{Interpretation:} The 8.88pp gap confirms spatial MLP reads a \emph{downstream echo} of the thermodynamic signal transformed through 7+ additional layers. The thermodynamic threshold accesses the \emph{primary output} of training---the delta magnitudes at layers 16--19 directly shaped by $\mathcal{L}_\text{thermo}$.

This is a general failure mode: gradient descent converges to whichever correlation is easiest to fit first. Spatial activations at layer 26 correlate with labels at $\sim$89\%, satisfying the loss. The gradient toward the underlying physics ($T(x)$) is negligibly small relative to the spatial shortcut, so the MLP never discovers it.

\subsection{Topology Metrics}

Table~\ref{tab:topology} shows cluster separation metrics before and after LoRA training (PCA-reduced activations, 2,000 examples).

\begin{table}[h]
\centering
\caption{Activation space topology metrics (t-SNE, UMAP, PCA).}
\label{tab:topology}
\begin{tabular}{lcccc}
\toprule
Method & Silhouette & Davies-Bouldin & Sep. Ratio & Inter-distance \\
\midrule
\multicolumn{5}{c}{\textit{Before LoRA training}} \\
PCA & 0.002 & 17.40 & 0.143 & 2.78 \\
\midrule
\multicolumn{5}{c}{\textit{After Phase 6 LoRA training}} \\
PCA & \textbf{0.226} & \textbf{1.58} & \textbf{1.547} & \textbf{31.1} \\
\midrule
Improvement & 113$\times$ & 11$\times$ better & 10.8$\times$ & 11.2$\times$ \\
\bottomrule
\end{tabular}
\end{table}

\textbf{Separation ratio $> 1$:} Clusters are further apart than they are wide, indicating genuine geometric structure rather than overlap.

\subsection{Context Sensitivity Control}

To verify the signal is epistemic (context-dependent) rather than phrasing-based, we run each example twice: with context and with context replaced by \texttt{[REMOVED]}.

\begin{table}[h]
\centering
\caption{Context sensitivity test (predictor score delta).}
\label{tab:context}
\begin{tabular}{lccc}
\toprule
Class & Score w/ context & Score w/o context & $\Delta$ \\
\midrule
Factual & 0.743 & 0.001 & --0.742 \\
Speculative & 0.153 & 0.002 & --0.151 \\
\bottomrule
\end{tabular}
\end{table}

\textbf{Sensitivity ratio:} 4.91$\times$ (factual is $\sim$5$\times$ more context-dependent). \textbf{Verdict:} Real epistemic signal, not phrasing artifact.

\section{Discussion}

\subsection{Why Thermodynamic Routing Outperforms Spatial Probes}

The 7--9pp performance gap between thermodynamic thresholds (98.39\%) and spatial MLPs (89--91\%) reveals a fundamental insight about mechanistic interpretability:

\textbf{Gradient descent optimizes for correlation, not causation.} The spatial MLP was given activations at layer 26, which correlate with epistemic state at $\sim$89\% accuracy. This correlation exists because layer 26 inherits structure from the thermodynamic process 7 layers earlier. The MLP converged to this \emph{echo} because it was sufficient to minimize BCE loss.

\textbf{The thermodynamic signal is the source.} Training directly shaped $\|\delta_n\|$ at layers 16--19 via $\mathcal{L}_\text{thermo}$. Reading $T(x)$ at these layers accesses the primary output of training with no downstream dilution.

This suggests a general principle: \textbf{interpretability methods should probe the direct outputs of training objectives, not downstream correlates.}

\subsection{Why Specialization Improves Generation}

The monotonic decrease in eval LM loss (1.914 $\rightarrow$ 1.835) as thermodynamic gap widened contradicts naive intuition that routing constraints would harm generation quality. We hypothesize:

\textbf{Uniform computation is a lazy equilibrium.} Cross-entropy loss does not penalize applying equal computational effort to all queries. A model that uses medium effort for everything and produces outputs good enough to minimize loss has no gradient signal to specialize.

\textbf{Thermodynamic constraint breaks lazy equilibrium.} Under $\mathcal{L}_\text{thermo}$, uniform computation incurs penalty. The model searches for alternatives and discovers that specialization \emph{reduces both routing and generation loss}:
\begin{itemize}
\item High-energy factual pathways $\rightarrow$ more reliable retrieval $\rightarrow$ better factual generation
\item Low-energy speculative pathways $\rightarrow$ smoother, less overconfident generation $\rightarrow$ better calibration
\end{itemize}

\textbf{Implication:} Thermodynamic routing may be part of the \emph{optimal generation strategy}, not a safety tax. This aligns with emerging evidence that mechanistic interpretability can guide capability improvements \citep{elhage2021mathematical}.

\subsection{Scaling Experiment: 1.5B to 7B}

Standard mechanistic interpretability degrades with scale: larger models have more feature superposition, more entangled representations, more resistance to spatial untangling \citep{elhage2022superposition}.

We directly tested whether thermodynamic routing scales by training the identical recipe on Qwen2.5-7B-Instruct (4.67$\times$ more parameters, same 28 layers, hidden dimension 1536$\rightarrow$3584) using two-GPU DDP via HuggingFace Accelerate. All thermodynamic hyperparameters were held constant.

\begin{table}[h]
\centering
\small
\begin{tabular}{lrrl}
\toprule
Metric & 1.5B & 7B & Change \\
\midrule
Factual mean $T(x)$ & 36.22 & 43.96 & $+21\%$ \\
Speculative mean $T(x)$ & 22.59 & 30.08 & $+33\%$ \\
\textbf{Gap} & \textbf{13.62} & \textbf{13.88} & $\mathbf{+1.9\%}$ \\
Best threshold $\tau$ & 27.2 & 35.4 & $+30\%$ \\
\textbf{Accuracy} & \textbf{98.39\%} & \textbf{98.13\%} & $\mathbf{-0.26\text{pp}}$ \\
\bottomrule
\end{tabular}
\caption{1.5B vs 7B thermodynamic evaluation (identical hyperparameters, same test set). Gap is essentially flat across a 4.67$\times$ parameter increase.}
\end{table}

The trained delta magnitude gap is essentially unchanged (13.62 $\rightarrow$ 13.88, implied scaling exponent $\alpha \approx 0.01$). Both models saturate at approximately $3.5\times$ the \texttt{thermo\_margin} (gap $\approx 13$--14 vs.\ margin $= 4.0$). This suggests the gap is controlled by the training objective, not model capacity: amplification saturates once the model has sufficient representational room to satisfy the constraint, which both models achieve within 6 epochs.

What does scale as expected is the absolute $T(x)$ level. The threshold shifts $+30\%$ (27.2$\rightarrow$35.4), broadly consistent with the larger residual stream norms ($\sqrt{3584/1536} = 1.53\times$ predicted; observed factual norm ratio $= 1.21\times$). Thresholds must be recalibrated per deployment; the \emph{gap} between classes is the reliable quantity.

\textbf{Interpretation:} The method transfers directly to a 4.67$\times$ larger model with no hyperparameter changes, and classification accuracy is essentially unchanged ($-0.26$\,pp). The thermodynamic routing approach is robust across this scale range. Whether the gap grows at much larger scales (70B, 405B) depends on whether the amplification budget is set by model capacity or solely by the loss margin and training duration; testing with a higher \texttt{thermo\_margin} would distinguish the two.

The anticorrelation between salience and curvature reported in \citet{manson2025curved} ($r = -0.89$ in LLaMA) yields an additional testable prediction: if factual processing corresponds to the high-$T(x)$ regime, it should exhibit lower curvature than speculative processing at the same layers. This is directly checkable on the existing 7B hidden states without retraining. More broadly, replacing the Euclidean norm with the $G$-weighted norm $\|h_{n+1} - h_n\|_G$ (where $G = U^\top U$) may tighten the distributional gap by projecting movement into token-probability space.

\subsection{Limitations and Caveats}

\textbf{Threshold calibration:} The absolute value $\tau = 27.2$ is model- and distribution-specific. The reliable quantity is the \emph{gap} between classes ($\sim$13.62), not absolute magnitudes. Thresholds must be calibrated on held-out data per deployment.

\textbf{Single model family:} Primary results are on Qwen2.5-1.5B; the 7B replication used the same family. Generalization across architectures (Llama, Mistral, Gemma) remains untested.

\textbf{Domain transfer:} Dataset is Wikipedia-based (SQuAD). Transfer to scientific papers, code, dialogue, or other domains requires empirical validation.

\textbf{Residual error:} 1.61\% test error (30 examples) not analyzed. Unclear whether genuinely ambiguous or correctable with better features/margins.

\textbf{Inference overhead:} Requires retaining hidden states at layers 16--19 during forward pass. Compute cost is negligible (norm operations on already-computed tensors) but requires hook access to mid-layer activations.

\section{Future Work}

\textbf{Cross-architecture validation:} Test whether thermodynamic signal exists in Llama, Mistral, GPT families without LoRA training (zero-shot transfer).

\textbf{Scale experiments:} The 7B replication (Section 5.4) showed the gap is hyperparameter-bounded rather than capacity-bounded in the 1.5B--7B range. Next: test whether increasing \texttt{thermo\_margin} proportionally increases the trained gap, and whether a qualitative change emerges at 70B.

\textbf{Domain generalization:} Train domain-specific predictors on medical (MedQA), legal (contracts), code (LeetCode), scientific (arXiv) corpora using identical architecture and threshold calibration procedure.

\textbf{Dynamic agentic contexts:} Track how $T(x)$ shifts in multi-turn conversations as tool calls add grounded information mid-generation, validating responsiveness for production deployment.

\textbf{Correlation with factual accuracy:} Test whether high $T(x)$ predicts correctness on open-ended generation (not just classification), potentially enabling real-time quality monitoring.

\textbf{Intervention strategies:} Develop production systems that trigger retrieval, chain-of-thought, or human review when $T(x)$ crosses calibrated thresholds.

\section{Conclusion}

We have demonstrated that large language models can be trained to develop thermodynamically distinct computational pathways for factual versus speculative generation, and that this distinction is more fundamentally readable via \emph{delta magnitudes} (layer-to-layer hidden state changes) than via spatial activation patterns. A zero-parameter threshold on $T(x) = \text{mean}(\|h_{n+1} - h_n\|)$ achieves 98.39\% classification accuracy on Qwen2.5-1.5B, outperforming trained spatial MLPs by 7--9pp on the same model. Direct replication on Qwen2.5-7B (4.67$\times$ parameters, identical hyperparameters, 2-GPU DDP) produced 98.13\% accuracy with the same gap magnitude, confirming the method transfers across model scale without modification.

This finding has implications for mechanistic interpretability: probes should target the \emph{direct outputs of training objectives}, not downstream correlates that gradient descent exploits as shortcuts. It also suggests that epistemic specialization may improve generation quality rather than harm it, contradicting typical capability-safety tradeoffs.

The approach requires no architectural changes, uses automated labeling, and produces interpretable structure suitable for production LLM systems. If the thermodynamic signal generalizes across model families and scales favorably with size, it offers a path toward mechanistically grounded uncertainty quantification at frontier scale---a critical requirement for deploying LLMs in high-stakes domains.

\section*{Code and Data Availability}

Code, trained models, and datasets are available at \texttt{https://github.com/ameritusweb/thermodynamic-epistemic-routing}.

\section*{License}

This work is licensed under the Apache License 2.0. The trained models and datasets are released under the same license for research purposes.

\bibliographystyle{plainnat}
\begin{thebibliography}{99}

\bibitem[Angelopoulos and Bates(2021)]{angelopoulos2021gentle}
Angelopoulos, A.~N. and Bates, S. (2021).
\newblock A gentle introduction to conformal prediction and distribution-free uncertainty quantification.
\newblock \emph{arXiv preprint arXiv:2107.07511}.

\bibitem[Bai et al.(2022)]{bai2022constitutional}
Bai, Y., Kadavath, S., Kundu, S., et al. (2022).
\newblock Constitutional AI: Harmlessness from AI feedback.
\newblock \emph{arXiv preprint arXiv:2212.08073}.

\bibitem[Belinkov(2021)]{belinkov2021probing}
Belinkov, Y. (2021).
\newblock Probing classifiers: Promises, shortcomings, and advances.
\newblock \emph{Computational Linguistics}, 48(1):207--219.

\bibitem[Blundell et al.(2015)]{blundell2015weight}
Blundell, C., Cornebise, J., Kavukcuoglu, K., and Wierstra, D. (2015).
\newblock Weight uncertainty in neural networks.
\newblock In \emph{ICML}, pages 1613--1622.

\bibitem[Brown et al.(2020)]{brown2020language}
Brown, T., Mann, B., Ryder, N., et al. (2020).
\newblock Language models are few-shot learners.
\newblock In \emph{NeurIPS}, volume 33, pages 1877--1901.

\bibitem[Cunningham et al.(2023)]{cunningham2023sparse}
Cunningham, H., Ewart, A., Riggs, L., Huben, R., and Sharkey, L. (2023).
\newblock Sparse autoencoders find highly interpretable features in language models.
\newblock \emph{arXiv preprint arXiv:2309.08600}.

\bibitem[Elhage et al.(2021)]{elhage2021mathematical}
Elhage, N., Hume, T., Olsson, C., et al. (2021).
\newblock A mathematical framework for transformer circuits.
\newblock \emph{Transformer Circuits Thread}.

\bibitem[Elhage et al.(2022)]{elhage2022superposition}
Elhage, N., Hume, T., Olsson, C., et al. (2022).
\newblock Toy models of superposition.
\newblock \emph{Transformer Circuits Thread}.

\bibitem[Gal and Ghahramani(2016)]{gal2016dropout}
Gal, Y. and Ghahramani, Z. (2016).
\newblock Dropout as a Bayesian approximation: Representing model uncertainty in deep learning.
\newblock In \emph{ICML}, pages 1050--1059.

\bibitem[Goodfellow et al.(2014)]{goodfellow2014generative}
Goodfellow, I., Pouget-Abadie, J., Mirza, M., et al. (2014).
\newblock Generative adversarial nets.
\newblock In \emph{NeurIPS}, pages 2672--2680.

\bibitem[Hu et al.(2021)]{hu2021lora}
Hu, E.~J., Shen, Y., Wallis, P., et al. (2021).
\newblock LoRA: Low-rank adaptation of large language models.
\newblock \emph{arXiv preprint arXiv:2106.09685}.

\bibitem[Irving et al.(2018)]{irving2018ai}
Irving, G., Christiano, P., and Amodei, D. (2018).
\newblock AI safety via debate.
\newblock \emph{arXiv preprint arXiv:1805.00899}.

\bibitem[Jiang et al.(2023)]{jiang2023mistral}
Jiang, A.~Q., Sablayrolles, A., Mensch, A., et al. (2023).
\newblock Mistral 7B.
\newblock \emph{arXiv preprint arXiv:2310.06825}.

\bibitem[Kuhn et al.(2023)]{kuhn2023semantic}
Kuhn, L., Gal, Y., and Farquhar, S. (2023).
\newblock Semantic uncertainty: Linguistic invariances for uncertainty estimation in natural language generation.
\newblock \emph{arXiv preprint arXiv:2302.09664}.

\bibitem[Lakshminarayanan et al.(2017)]{lakshminarayanan2017simple}
Lakshminarayanan, B., Pritzel, A., and Blundell, C. (2017).
\newblock Simple and scalable predictive uncertainty estimation using deep ensembles.
\newblock In \emph{NeurIPS}, pages 6402--6413.

\bibitem[Lewis et al.(2020)]{lewis2020retrieval}
Lewis, P., Perez, E., Piktus, A., et al. (2020).
\newblock Retrieval-augmented generation for knowledge-intensive NLP tasks.
\newblock In \emph{NeurIPS}, volume 33, pages 9459--9474.

\bibitem[Lialin et al.(2023)]{lialin2023scaling}
Lialin, V., Deshpande, V., and Rumshisky, A. (2023).
\newblock Scaling down to scale up: A guide to parameter-efficient fine-tuning.
\newblock \emph{arXiv preprint arXiv:2303.15647}.

\bibitem[Manson(2025)]{manson2025curved}
Manson, R. (2025).
\newblock Curved inference: Concern-sensitive geometry in large language model residual streams.
\newblock \emph{arXiv preprint arXiv:2507.21107}.

\bibitem[Meng et al.(2022)]{meng2022locating}
Meng, K., Bau, D., Andonian, A., and Belinkov, Y. (2022).
\newblock Locating and editing factual associations in GPT.
\newblock In \emph{NeurIPS}, volume 35, pages 17359--17372.

\bibitem[Qwen Team(2024)]{qwen2024}
Qwen Team (2024).
\newblock Qwen2.5: A party of foundation models.
\newblock \emph{https://qwenlm.github.io/blog/qwen2.5/}.

\bibitem[Rajpurkar et al.(2018)]{rajpurkar2018know}
Rajpurkar, P., Jia, R., and Liang, P. (2018).
\newblock Know what you don't know: Unanswerable questions for SQuAD.
\newblock In \emph{ACL}, pages 784--789.

\bibitem[Tian et al.(2023)]{tian2023finetuning}
Tian, K., Mitchell, E., Yao, H., Manning, C.~D., and Finn, C. (2023).
\newblock Fine-tuning language models for factuality.
\newblock \emph{arXiv preprint arXiv:2311.08401}.

\bibitem[Touvron et al.(2023)]{touvron2023llama}
Touvron, H., Martin, L., Stone, K., et al. (2023).
\newblock Llama 2: Open foundation and fine-tuned chat models.
\newblock \emph{arXiv preprint arXiv:2307.09288}.

\bibitem[Wang et al.(2023)]{wang2023selfconsistency}
Wang, X., Wei, J., Schuurmans, D., et al. (2023).
\newblock Self-consistency improves chain of thought reasoning in language models.
\newblock In \emph{ICLR}.

\end{thebibliography}

\end{document}
